
\documentclass[ukrainian]{ucaphyslab}

\labtitle{Вивчення рівноприскореного руху за допомогою машини Атвуда}
\labsubtitle{Лабораторна робота № 4}

\student{Довбенчук Олена}
\student{Кульбаба Андрій}
\student{Магарита Ія}
\student{Пастарнак Марко}

\makeatletter
\renewcommand{\ucu@labelstudent}{Студенти}
\renewcommand{\ucu@labelteacher}{Викладач}
\makeatother

\teacher{Володимир Коломієць}

\labdate{02-10-2026}

\usepackage{float}
\usepackage{pgfplots}
\pgfplotsset{compat=1.18}
\sisetup{uncertainty-mode=separate}
\definecolor{series1}{HTML}{2673d1}
\definecolor{series2}{HTML}{EB6834}

\begin{document}
\maketitlepage

\section{Мета роботи}
\begin{enumerate}
  \item Визначити прискорення вільного падіння за допомогою машини Атвуда.
  \item Вивчити основні закономірності прямолінійного рівноприскореного руху.
\end{enumerate}

\section{Прилади та обладнання}
\begin{itemize}
  \item Машина Атвуда: блок, нитка, два важки однакової маси $M$;
  \item набір додаткових тягарців;
  \item лінійка (ціна поділки \SI{1}{\milli\meter});
  \item секундоміри (час вимірювали 4 спостерігачі одночасно).
\end{itemize}

\section{Теоретичні відомості}
Знехтуємо опором повітря, тертям у блоці, масою нитки та блока. Тоді сили натягу нитки
з обох боків однакові, а рівняння руху важків у проєкції на вісь $x$:
\begin{align}
  Ma &= T - Mg, \tag{1.13}\\
  (M+m)a &= (M+m)g - T, \tag{1.14}
\end{align}
звідки
\begin{equation}
  a = g\,\frac{m}{m+2M}. \tag{1.15}
\end{equation}
Прискорення стале, тобто рух рівноприскорений. Час, за який важок проходить відстань $h$:
\begin{equation}
  t = \sqrt{\frac{2h}{a}} = \sqrt{\frac{2h}{g}\cdot\frac{m+2M}{m}}. \tag{1.16}
\end{equation}

З (1.16) прискорення важків $a = 2h/t^2$, а прискорення вільного падіння
\begin{equation}
  g = a\,\frac{m+2M}{m}. \tag{1.17}
\end{equation}
Якщо врахувати тертя в блоці (масу блока не враховуємо, $I=0$), маса $m_0$, за якої
починається рух, задає момент сил тертя: $M_\text{т}/r = m_0 g$. Тоді
\begin{equation}
  g = \frac{(m+2M)\,a}{m - m_0}. \tag{1.26}
\end{equation}

Якщо рух рівноприскорений, то за (1.16) при $m=\text{const}$ залежність $t(\sqrt h)$ лінійна,
а при $h=\text{const}$ і $m \ll 2M$ лінійна залежність $t\bigl(\sqrt{M/m}\bigr)$:
$t \approx 2\sqrt{h/g}\,\sqrt{M/m}$.

\section{Хід роботи та результати вимірювань}
\begin{enumerate}
  \item Визначили масу $m_0$ додаткового тягарця, за якої починається рух системи.
  \item Для фіксованої маси $m$ виміряли час руху важка для 5 висот $h$.
  \item Для фіксованої висоти $h_\text{max}$ виміряли час руху для 5 мас $m$.
\end{enumerate}
Кожну спробу проводили 5 разів. Час кожної спроби одночасно вимірювали 4 спостерігачі;
у таблицях $t_i$~— середнє за спостерігачами, первинні дані наведено в додатку.

\begin{table}[H]
  \centering
  \begin{tabular}{@{}llc@{}}
    \toprule
    Величина & Позначення & Значення \\
    \midrule
    Маса основного важка & $M$ & \SI{106.50}{\gram} \\
    Маса, за якої починається рух & $m_0$ & \SI{5.14}{\gram} \\
    Маса тягарця (експеримент 1) & $m$ & \SI{51.17}{\gram} \\
    Висота (експеримент 2) & $h_\text{max}$ & \SI{0.975}{\meter} \\
    Табличне значення & $g$ & \SI{9.81}{\meter\per\second\squared} \\
    \bottomrule
  \end{tabular}
  \caption{Константи та параметри установки}
\end{table}

Похибка часу: $\Delta t = t_{0.95;\,4}\,\sqrt{\dfrac{\sum(t_i-\bar t)^2}{n(n-1)}}$,
$t_{0.95;\,4}=2.78$, $n=5$. Похибка висоти $\Delta h = \SI{1}{\milli\meter}$ (ціна поділки лінійки).

\begin{table}[H]
  \centering
  \begin{tabular}{@{}c ccccc ccc@{}}
    \toprule
    $h$, м & $t_1$, с & $t_2$, с & $t_3$, с & $t_4$, с & $t_5$, с & $\bar t$, с & $\Delta t$, с & $\Delta h$, м \\
    \midrule
    0.975 & 1.260 & 0.965 & 1.123 & 1.028 & 1.023 & 1.080 & 0.144 & 0.001 \\
    0.918 & 0.958 & 0.880 & 1.003 & 0.988 & 0.983 & 0.962 & 0.060 & 0.001 \\
    0.863 & 0.913 & 0.940 & 1.033 & 1.093 & 0.960 & 0.988 & 0.091 & 0.001 \\
    0.802 & 0.900 & 0.938 & 0.843 & 0.895 & 0.850 & 0.885 & 0.049 & 0.001 \\
    0.746 & 0.888 & 0.833 & 0.823 & 0.828 & 0.845 & 0.843 & 0.033 & 0.001 \\
    \bottomrule
  \end{tabular}
  \caption{Залежність часу руху від висоти ($m = \SI{51.17}{\gram}$)}
\end{table}

\begin{table}[H]
  \centering
  \begin{tabular}{@{}cc ccccc cc@{}}
    \toprule
    $m$, г & $\sqrt{M/m}$ & $t_1$, с & $t_2$, с & $t_3$, с & $t_4$, с & $t_5$, с & $\bar t$, с & $\Delta t$, с \\
    \midrule
    30.32 & 1.874 & 1.478 & 1.365 & 1.475 & 1.443 & 1.345 & 1.421 & 0.077 \\
    35.60 & 1.730 & 1.298 & 1.283 & 1.350 & 1.303 & 1.255 & 1.298 & 0.043 \\
    40.67 & 1.618 & 1.265 & 1.125 & 1.170 & 1.150 & 1.250 & 1.192 & 0.077 \\
    45.97 & 1.522 & 1.095 & 1.105 & 1.030 & 1.148 & 1.040 & 1.084 & 0.060 \\
    51.17 & 1.443 & 0.963 & 0.943 & 1.033 & 0.898 & 1.005 & 0.968 & 0.066 \\
    \bottomrule
  \end{tabular}
  \caption{Залежність часу руху від маси тягарця ($h = \SI{0.975}{\meter}$)}
\end{table}

\section{Обробка результатів}

\subsection{Експеримент 1: залежність від висоти}

\begin{figure}[H]
  \centering
  \begin{tikzpicture}
    \begin{axis}[
      width=0.85\linewidth, height=6.5cm,
      xlabel={$\sqrt{h}$, м$^{1/2}$}, ylabel={$\bar t$, с},
      xmin=0.85, xmax=1.0, ymin=0.75, ymax=1.25,
      ymajorgrids, grid style={gray!25},
      axis lines*=left,
      x tick label style={/pgf/number format/.cd, fixed, precision=2},
      legend style={draw=none, at={(0.5,1.02)}, anchor=south, legend columns=2, font=\small},
    ]
      \addplot[color=series1, only marks, mark=*, mark size=2pt,
        error bars/.cd, y dir=both, y explicit] coordinates {
        (0.9874,1.080) +- (0,0.144) (0.9581,0.962) +- (0,0.060) (0.9290,0.988) +- (0,0.091)
        (0.8955,0.885) +- (0,0.049) (0.8637,0.843) +- (0,0.033)};
    \end{axis}
  \end{tikzpicture}
  \caption{Залежність часу руху від $\sqrt h$}
\end{figure}

Прискорення та $g$ для кожної висоти окремо ($a = 2h/\bar t^{\,2}$, $g$ за (1.17)):

\begin{table}[H]
  \centering
  \begin{tabular}{@{}ccccc@{}}
    \toprule
    $h$, м & $\bar t$, с & $\bar t^{\,2}$, с$^2$ & $a$, м/с$^2$ & $g$, м/с$^2$ \\
    \midrule
    0.975 & 1.080 & 1.165 & 1.673 & 8.64 \\
    0.918 & 0.962 & 0.925 & 1.984 & 10.24 \\
    0.863 & 0.988 & 0.975 & 1.770 & 9.14 \\
    0.802 & 0.885 & 0.783 & 2.048 & 10.57 \\
    0.746 & 0.843 & 0.711 & 2.099 & 10.84 \\
    \midrule
    Сер. & & & 1.915 & 9.89 \\
    \bottomrule
  \end{tabular}
  \caption{Прискорення $a$ та $g$ (експеримент 1)}
\end{table}

\subsection{Експеримент 2: залежність від маси}

\begin{figure}[H]
  \centering
  \begin{tikzpicture}
    \begin{axis}[
      width=0.85\linewidth, height=6.5cm,
      xlabel={$\sqrt{M/m}$}, ylabel={$\bar t$, с},
      xmin=1.4, xmax=1.9, ymin=0.85, ymax=1.55,
      ymajorgrids, grid style={gray!25},
      axis lines*=left,
      legend style={draw=none, at={(0.5,1.02)}, anchor=south, legend columns=2, font=\small},
    ]
      \addplot[color=series1, only marks, mark=*, mark size=2pt,
        error bars/.cd, y dir=both, y explicit] coordinates {
        (1.8742,1.421) +- (0,0.077) (1.7296,1.298) +- (0,0.043) (1.6182,1.192) +- (0,0.077)
        (1.5221,1.084) +- (0,0.060) (1.4427,0.968) +- (0,0.066)};
    \end{axis}
  \end{tikzpicture}
  \caption{Залежність часу руху від $\sqrt{M/m}$}
\end{figure}

Прискорення та $g$ для кожної маси окремо ($a = 2h/\bar t^{\,2}$, $g$ за (1.17)):

\begin{table}[H]
  \centering
  \begin{tabular}{@{}cccc@{}}
    \toprule
    $m$, г & $\bar t$, с & $a$, м/с$^2$ & $g$, м/с$^2$ \\
    \midrule
    30.32 & 1.421 & 0.966 & 7.75 \\
    35.60 & 1.298 & 1.158 & 8.09 \\
    40.67 & 1.192 & 1.372 & 8.56 \\
    45.97 & 1.084 & 1.661 & 9.36 \\
    51.17 & 0.968 & 2.081 & 10.74 \\
    \midrule
    Сер. & & 1.448 & 8.90 \\
    \bottomrule
  \end{tabular}
  \caption{Прискорення $a$ та $g$ (експеримент 2)}
\end{table}

\subsection{Порівняння з табличним значенням}

\begin{table}[H]
  \centering
  \begin{tabular}{@{}lc@{}}
    \toprule
    Спосіб & $g$, м/с$^2$ \\
    \midrule
    Експеримент 1, середнє за висотами (Табл.~2.1) & \num{9.89} \\
    Експеримент 2, середнє за масами (Табл.~2.2) & \num{8.90} \\
    \midrule
    Табличне значення & \num{9.81} \\
    \bottomrule
  \end{tabular}
  \caption{Значення прискорення вільного падіння}
\end{table}

\section{Висновки}
У ході лабораторної роботи було досліджено рівноприскорений рух за допомогою машини Атвуда.
Прискорення вільного падіння $g$ визначено двома способами: змінюючи висоту падіння та
змінюючи масу тягарця. Отримані результати порівняно з табличним значенням
$g = \SI{9.81}{\meter\per\second\squared}$.

Експериментальні залежності $t(\sqrt h)$ та $t\bigl(\sqrt{M/m}\bigr)$ у межах похибок є лінійними,
що підтверджує рівноприскорений характер руху. За першим способом середнє значення
$g = \SI{9.89}{\meter\per\second\squared}$ (відхилення від табличного $0.8\,\%$), за другим~—
$g = \SI{8.90}{\meter\per\second\squared}$ (відхилення $9.3\,\%$). Отже, відносне відхилення
не перевищує $10\,\%$, що є прийнятним для даного експерименту.

Основні причини відхилень: ручне вимірювання секундоміром короткого (близько \SI{1}{\second})
часу руху.
\section*{Додаток. Первинні дані}

\begin{table}[H]
  \centering
  \small
  \begin{tabular}{@{}cc cccc c@{}}
    \toprule
    $h$, м & Спроба & Сп.~1, с & Сп.~2, с & Сп.~3, с & Сп.~4, с & Сер., с \\
    \midrule
    0.975 & 1 & 1.06 & 1.06 & 1.44 & 1.48 & 1.260 \\
     & 2 & 0.81 & 0.90 & 0.98 & 1.17 & 0.965 \\
     & 3 & 1.13 & 0.98 & 1.21 & 1.17 & 1.123 \\
     & 4 & 1.19 & 0.73 & 0.98 & 1.21 & 1.028 \\
     & 5 & 0.87 & 1.30 & 0.98 & 0.94 & 1.023 \\
    \midrule
    0.918 & 1 & 0.86 & 1.19 & 0.98 & 0.80 & 0.958 \\
     & 2 & 0.81 & 1.02 & 0.86 & 0.83 & 0.880 \\
     & 3 & 0.94 & 1.18 & 1.03 & 0.86 & 1.003 \\
     & 4 & 1.00 & 1.10 & 0.94 & 0.91 & 0.988 \\
     & 5 & 1.00 & 1.17 & 0.93 & 0.83 & 0.983 \\
    \midrule
    0.863 & 1 & 0.81 & 0.81 & 1.08 & 0.95 & 0.913 \\
     & 2 & 0.81 & 0.83 & 0.94 & 1.18 & 0.940 \\
     & 3 & 0.80 & 1.26 & 1.14 & 0.93 & 1.033 \\
     & 4 & 0.92 & 1.30 & 0.98 & 1.17 & 1.093 \\
     & 5 & 0.87 & 0.99 & 1.17 & 0.81 & 0.960 \\
    \midrule
    0.802 & 1 & 0.93 & 0.98 & 0.79 & 0.90 & 0.900 \\
     & 2 & 0.67 & 1.01 & 1.06 & 1.01 & 0.938 \\
     & 3 & 0.74 & 0.84 & 0.84 & 0.95 & 0.843 \\
     & 4 & 0.80 & 1.01 & 0.94 & 0.83 & 0.895 \\
     & 5 & 0.73 & 0.98 & 0.83 & 0.86 & 0.850 \\
    \midrule
    0.746 & 1 & 0.81 & 1.02 & 0.91 & 0.81 & 0.888 \\
     & 2 & 0.67 & 1.05 & 0.78 & 0.83 & 0.833 \\
     & 3 & 0.73 & 0.92 & 0.83 & 0.81 & 0.823 \\
     & 4 & 0.67 & 0.88 & 0.76 & 1.00 & 0.828 \\
     & 5 & 0.80 & 0.93 & 0.97 & 0.68 & 0.845 \\
    \bottomrule
  \end{tabular}
  \caption{Час руху для різних висот за спостерігачами ($m = \SI{51.17}{\gram}$)}
\end{table}

\begin{table}[H]
  \centering
  \small
  \begin{tabular}{@{}cc cccc c@{}}
    \toprule
    $m$, г & Спроба & Сп.~1, с & Сп.~2, с & Сп.~3, с & Сп.~4, с & Сер., с \\
    \midrule
    30.32 & 1 & 1.41 & 1.46 & 1.46 & 1.58 & 1.478 \\
     & 2 & 1.20 & 1.45 & 1.41 & 1.40 & 1.365 \\
     & 3 & 1.46 & 1.49 & 1.49 & 1.46 & 1.475 \\
     & 4 & 1.27 & 1.38 & 1.42 & 1.70 & 1.443 \\
     & 5 & 1.26 & 1.43 & 1.44 & 1.25 & 1.345 \\
    \midrule
    35.60 & 1 & 1.19 & 1.26 & 1.31 & 1.43 & 1.298 \\
     & 2 & 1.20 & 1.26 & 1.44 & 1.23 & 1.283 \\
     & 3 & 1.34 & 1.36 & 1.39 & 1.31 & 1.350 \\
     & 4 & 1.14 & 1.40 & 1.44 & 1.23 & 1.303 \\
     & 5 & 1.14 & 1.48 & 1.35 & 1.05 & 1.255 \\
    \midrule
    40.67 & 1 & 1.33 & 1.37 & 1.26 & 1.10 & 1.265 \\
     & 2 & 1.00 & 1.13 & 1.29 & 1.08 & 1.125 \\
     & 3 & 1.13 & 1.14 & 1.31 & 1.10 & 1.170 \\
     & 4 & 1.01 & 0.98 & 1.37 & 1.24 & 1.150 \\
     & 5 & 1.19 & 1.30 & 1.29 & 1.22 & 1.250 \\
    \midrule
    45.97 & 1 & 1.07 & 0.99 & 1.12 & 1.20 & 1.095 \\
     & 2 & 0.99 & 1.20 & 1.23 & 1.00 & 1.105 \\
     & 3 & 1.06 & 1.20 & 1.03 & 0.83 & 1.030 \\
     & 4 & 1.08 & 1.20 & 1.21 & 1.10 & 1.148 \\
     & 5 & 1.00 & 1.25 & 1.11 & 0.80 & 1.040 \\
    \midrule
    51.17 & 1 & 0.81 & 1.07 & 1.06 & 0.91 & 0.963 \\
     & 2 & 0.93 & 1.03 & 0.90 & 0.91 & 0.943 \\
     & 3 & 1.00 & 0.88 & 1.16 & 1.09 & 1.033 \\
     & 4 & 0.73 & 0.81 & 0.98 & 1.07 & 0.898 \\
     & 5 & 0.93 & 1.13 & 1.01 & 0.95 & 1.005 \\
    \bottomrule
  \end{tabular}
  \caption{Час руху для різних мас за спостерігачами ($h = \SI{0.975}{\meter}$)}
\end{table}

\end{document}
